% Môn: Toán
% Lớp: 11
% Chương: Chương I. Hàm số lượng giác và phương trình lượng giác
% Bài: Bài 1. Giá trị lượng giác của góc lượng giác
% Số câu: 8
% App đọc file này trực tiếp — sửa rồi Commit trên GitHub, bấm Đồng bộ GitHub .tex

% Bài 1. Giá trị lượng giác của góc lượng giác

% ===== Câu 1 =====
% ID: T11C1B1-01-DS
% Mức: TH
\dangbt{Chuyển đổi số đo góc giữa độ và radian}
\begin{ex}
% ID: T11C1B1-01-DS
% Mức: TH
Đổi số đo của các góc sang độ. Khẳng định sau đây \textbf{ĐÚNG} hay $\textbf{SAI}$ ?
\choiceTF

\loigiai{
A. Đúng — Ta có $\dfrac{3\pi}{4} \text{ rad} = \dfrac{3\pi}{4} \cdot \dfrac{180^\circ}{\pi} = 3 \cdot 45^\circ = 135^\circ$.
B. Đúng — Ta có $-\dfrac{\pi}{360} \text{ rad} = -\dfrac{\pi}{360} \cdot \dfrac{180^\circ}{\pi} = -\dfrac{1}{2}^\circ = -0.5^\circ$.
C. Sai — Ta có $\dfrac{31\pi}{2} \text{ rad} = \dfrac{31\pi}{2} \cdot \dfrac{180^\circ}{\pi} = 31 \cdot 90^\circ = 2790^\circ$. Mệnh đề cho là $27^\circ$, do đó sai.
D. Đúng — Ta có $-4 \text{ rad} = -4 \cdot \dfrac{180^\circ}{\pi} = -\dfrac{720^\circ}{\pi} \approx -229.18^\circ$.
}
\end{ex}

% ===== Câu 2 =====
% ID: T11C1B1-02-DS
% Mức: TH
\begin{ex}
% ID: T11C1B1-02-DS
% Mức: TH
Đổi số đo của các góc sang radian. Khi đó các mệnh đề sau \textbf{ĐÚNG} hay $\textbf{SAI}$?
\choiceTF

\loigiai{
A. Đúng — Ta có $30^{\circ} = 30 \times \dfrac{\pi}{180} = \dfrac{\pi}{6}$ rad.
B. Đúng — Ta có $\left(\dfrac{15}{\pi}\right)^{\circ} = \dfrac{15}{\pi} \times \dfrac{\pi}{180} = \dfrac{15}{180} = \dfrac{1}{12}$ rad.
C. Đúng — Ta có $132^{\circ} = 132 \times \dfrac{\pi}{180} = \dfrac{132}{180}\pi = \dfrac{11}{15}\pi$ rad.
D. Sai — Ta có $-495^{\circ} = -495 \times \dfrac{\pi}{180} = -\dfrac{495}{180}\pi = -\dfrac{11}{4}\pi$ rad. Mệnh đề cho là $-\dfrac{13\pi}{4}$ rad.
}
\end{ex}

% ===== Câu 3 =====
% ID: T11C1B1-03-DS
% Mức: TH
\dangbt{Xác định góc có cùng điểm biểu diễn trên đường tròn}
\begin{ex}
% ID: T11C1B1-03-DS
% Mức: TH
Các góc có cùng điểm biểu diễn trên đường tròn lượng giác \textbf{ĐÚNG} hay $\textbf{SAI}$?
\choiceTF

\loigiai{
A. Đúng — Hiệu của hai góc là $756^{\circ} - (-324^{\circ}) = 756^{\circ} + 324^{\circ} = 1080^{\circ} = 3 \times 360^{\circ}$.
B. Đúng — Hiệu của hai góc là $-324^{\circ} - 36^{\circ} = -360^{\circ} = -1 \times 360^{\circ}$.
C. Sai — Hiệu của hai góc là $216^{\circ} - 36^{\circ} = 180^{\circ}$, không phải là bội số nguyên của $360^{\circ}$.
D. Đúng — Hiệu của hai góc là $-\dfrac{41\pi}{7} - \dfrac{15\pi}{7} = -\dfrac{56\pi}{7} = -8\pi = -4 \times 2\pi$.
}
\end{ex}

% ===== Câu 4 =====
% ID: T11C1B1-04-DS
% Mức: TH
\begin{ex}
% ID: T11C1B1-04-DS
% Mức: TH
Các cặp góc sau có cùng điểm biểu diễn trên đường tròn lượng giác \textbf{ĐÚNG} hay $\textbf{SAI}$ ?
\choiceTF

\loigiai{
A. Đúng — Ta có $1127^\circ = 3 \times 360^\circ + 47^\circ$ và $-313^\circ = -1 \times 360^\circ + 47^\circ$. Vì $1127^\circ - (-313^\circ) = 1440^\circ = 4 \times 360^\circ$, hai góc này có cùng điểm biểu diễn.
B. Sai — Ta có $1127^\circ = 3 \times 360^\circ + 47^\circ$ và $-674^\circ = -2 \times 360^\circ + 46^\circ$. Vì $1127^\circ - (-674^\circ) = 1801^\circ$ không phải là bội số của $360^\circ$, hai góc này không có cùng điểm biểu diễn.
C. Đúng — Ta có $\dfrac{61\pi}{5} = 12\pi + \dfrac{\pi}{5} = 6 \times 2\pi + \dfrac{\pi}{5}$ và $-\dfrac{19\pi}{5} = -4\pi + \dfrac{\pi}{5} = -2 \times 2\pi + \dfrac{\pi}{5}$. Vì $\dfrac{61\pi}{5} - \left(-\dfrac{19\pi}{5}\right) = \dfrac{80\pi}{5} = 16\pi = 8 \times 2\pi$, hai góc này có cùng điểm biểu diễn.
D. Sai — Ta có $\dfrac{61\pi}{5} = 12\pi + \dfrac{\pi}{5}$ và $-\dfrac{23\pi}{4} = -6\pi - \dfrac{\pi}{4}$. Vì $\dfrac{61\pi}{5} - \left(-\dfrac{23\pi}{4}\right) = \dfrac{244\pi + 115\pi}{20} = \dfrac{359\pi}{20}$ không phải là bội số của $2\pi$, hai góc này không có cùng điểm biểu diễn.
}
\end{ex}

% ===== Câu 5 =====
% ID: T11C1B1-05-DS
% Mức: VDC
\begin{ex}
% ID: T11C1B1-05-DS
% Mức: VDC
$\immini$ {Trong hình vẽ bên, ta xem hình ảnh đường tròn trên một bánh lái tàu thuỷ tương ứng với một đường tròn lượng giác. Các mệnh đề sau \textbf{ĐÚNG} hay $\textbf{SAI}$ ?
\choiceTF

\loigiai{
A. Sai — Từ hình vẽ, $OA$ ứng với góc $0$ và $OB$ ứng với góc $\dfrac{\pi}{4}$, nên góc lượng giác $(OA, OB) = \dfrac{\pi}{4} + k2\pi\,(k \in \mathbb{Z})$. Tuy nhiên, đề bài ghi nhận $OB$ ứng với góc $\dfrac{\pi}{2}$, suy ra $(OA, OB) = \dfrac{\pi}{2} + k2\pi\,(k \in \mathbb{Z})$, không phải $\dfrac{\pi}{4} + k2\pi$. Vậy mệnh đề $A$ sai.
B. Sai — Bốn điểm $A, C, E, G$ lần lượt ứng với các góc $0$, $\dfrac{\pi}{2}$, $\pi$, $\dfrac{3\pi}{2}$. Khoảng cách giữa hai góc liên tiếp là $\dfrac{\pi}{2}$, nên công thức tổng quát là $k\dfrac{\pi}{2}\,(k \in \mathbb{Z})$, không phải $k\dfrac{\pi}{3}$. Vậy mệnh đề $B$ sai.
C. Đúng — iểm $A$ ứng với góc $0^\circ$ (tức $k \cdot 360^\circ$) và điểm $E$ ứng với góc $180^\circ$ (tức $180^\circ + k \cdot 360^\circ$). Gộp lại, các góc lượng giác biểu diễn cho hai điểm $A$ và $E$ đều có dạng $k \cdot 180^\circ\,(k \in \mathbb{Z})$. Vậy mệnh đề $C$ đúng.
D. Sai — Theo hệ thức Saclơ, $(OA, OC) + (OC, OH) = (OA, OH)$. Từ hình vẽ, $OC$ ứng với $\dfrac{\pi}{2}$ và $OH$ ứng với $\dfrac{7\pi}{4}$, nên $(OA, OC) = \dfrac{\pi}{2}$ và $(OC, OH) = \dfrac{7\pi}{4} - \dfrac{\pi}{2} = \dfrac{5\pi}{4}$. Do đó $(OA, OH) = \dfrac{\pi}{2} + \dfrac{5\pi}{4} = \dfrac{7\pi}{4}$, công thức tổng quát là $\dfrac{7\pi}{4} + k2\pi\,(k \in \mathbb{Z})$, không phải $\dfrac{\pi}{4} + k2\pi$. Vậy mệnh đề $D$ sai.
}
\end{ex}

% ===== Câu 6 =====
% ID: T11C1B1-06-DS
% Mức: TH
\dangbt{Xác định góc phần tư của góc lượng giác}
\begin{ex}
% ID: T11C1B1-06-DS
% Mức: TH
Các mệnh đề sau đây \textbf{ĐÚNG} hay $\textbf{SAI}$?
\choiceTF

\loigiai{
A. Đúng — Điểm biểu diễn của góc lượng giác có số đo $218^\circ$ nằm ở góc phần tư thứ III, vì $180^\circ \lt  218^\circ \lt  270^\circ$.
B. Sai — Điểm biểu diễn của góc lượng giác có số đo $-405^\circ$ tương đương với góc $-45^\circ$, nằm ở góc phần tư thứ IV, nhưng mệnh đề sai ở chỗ $\widehat{AON} = -45^\circ$, không phải $45^\circ$.
C. Đúng — Điểm biểu diễn của góc lượng giác có số đo $\dfrac{25\pi}{4}$ tương đương với góc $\dfrac{\pi}{4}$, nằm ở góc phần tư thứ $I$.
D. Sai — Điểm biểu diễn của góc lượng giác có số đo $\dfrac{15\pi}{2}$ tương đương với góc $-\dfrac{\pi}{2}$, nhưng điểm $Q(0;-1)$ không thoả mãn điều kiện này vì $\widehat{AOQ} = -\dfrac{\pi}{2}$, không phải $\dfrac{\pi}{2}$.
}
\end{ex}

% ===== Câu 7 =====
% ID: T11C1B1-07-DS
% Mức: VDC
\begin{ex}
% ID: T11C1B1-07-DS
% Mức: VDC
Biểu diễn góc lượng giác trên đường tròn lượng giác. Các mệnh đề sau \textbf{ĐÚNG} hay $\textbf{SAI}$ ?
\choiceTF
{$36^{\circ}+k360^{\circ}, \, k\in \mathbb{Z}$ là điểm $M$ thuộc góc phần tư thứ $II$}
{\True $-60^{\circ}+k180^{\circ},\, k\in \mathbb{Z}$ là các điểm $M_1$, $M_2$ thuộc góc phần tư thứ $II$ và $IV$}
{$-\dfrac{\pi}{4}+k2\pi$, $k\in \mathbb{Z}$ là $M$ thuộc góc phần tư thứ $III$}
{\True $-\dfrac{\pi}{6}+k\dfrac{\pi}{2}$, $k\in \mathbb{Z}$ là bốn điểm $M$, $N$, $P$, $Q$ thuộc góc phần tư thứ $I$, $II$, $III$, $IV$}
\loigiai{
\begin{itemchoice}
\item (Sai) $36^{\circ}+k360^{\circ}, \, k\in \mathbb{Z}$ là điểm $M$ thuộc góc phần tư thứ $II$. (Vì): Góc $36^\circ + k \cdot 360^\circ$, với $k \in \mathbb{Z}$, có điểm biểu diễn thuộc góc phần tư thứ $I$, không phải góc phần tư thứ II
\item (Đúng) $-60^{\circ}+k180^{\circ},\, k\in \mathbb{Z}$ là các điểm $M_1$, $M_2$ thuộc góc phần tư thứ $II$ và $IV$. (Vì): Góc $-60^\circ + k \cdot 180^\circ$, với $k \in \mathbb{Z}$, có các điểm biểu diễn thuộc góc phần tư thứ II và IV
\item (Sai) $-\dfrac{\pi}{4}+k2\pi$, $k\in \mathbb{Z}$ là $M$ thuộc góc phần tư thứ $III$. (Vì): Góc $-\dfrac{\pi}{4} + k \cdot 2\pi$, với $k \in \mathbb{Z}$, có điểm biểu diễn thuộc góc phần tư thứ IV, không phải góc phần tư thứ III
\item (Đúng) $-\dfrac{\pi}{6}+k\dfrac{\pi}{2}$, $k\in \mathbb{Z}$ là bốn điểm $M$, $N$, $P$, $Q$ thuộc góc phần tư thứ $I$, $II$, $III$, $IV$. (Vì): Góc $-\dfrac{\pi}{6} + k \cdot \dfrac{\pi}{2}$, với $k \in \mathbb{Z}$, có bốn điểm biểu diễn thuộc góc phần tư thứ $I$, II, III, IV
\end{itemchoice}
}
\end{ex}

% ===== Câu 8 =====
% ID: T11C1B1-07-TN
% Mức: NB
\dangbt{Tính độ dài cung tròn khi biết bán kính và số đo góc}
\begin{ex}
% ID: T11C1B1-07-TN
% Mức: NB
Cho đường tròn có bán kính bằng $9$ (cm). Tìm số đo (theo radian) của cung có độ dài $3\pi$ (cm).
\choice
{\True $\dfrac{\pi}{3}$}
{$\dfrac{2\pi}{3}$}
{$\dfrac{\pi}{4}$}
{$\dfrac{\pi}{6}$}
\loigiai{
$l=R\alpha \Rightarrow \alpha=\dfrac{l}{R}=\dfrac{3\pi}{9}=\dfrac{\pi}{3}.$
}
\end{ex}
